\subsection{Experiment 4G: independent-fleet and trajectory validation}
\label{sec:experiment4g}
\paragraph{Purpose and frozen protocol.}
Experiment 4F recovered the expected architecture advantage on one fleet, but
used the same clients for identification and testing. Experiment 4G tests whether
that recovery transfers to independently sampled vehicles and additional
trajectories. Training seeds are 11--20, paired with test seeds 111--120. Each
fleet contains ten clients in each of the three unchanged physical families.
Thus 300 training and 300 separate test vehicles are used. Test clients share
the family distribution, not the realized physical parameters, with training
clients. Repeated textual client IDs are scoped by fleet seed.

All estimator bounds, starts, residual scaling, geometry assumptions, controller
weights, and speed grids remain those of 4F. The physics controller used for
data collection is constructed from the training fleet only. Structured fitting
uses near-linear training records only. No test trajectories or test-plant
matrices enter identification, tuning, or candidate selection. True test plants
are accessed only for simulation and the post-design frozen audit. Family labels
remain oracle information. The JSON protocol is frozen before the runs and its
SHA-256 digest is stored with the conclusions.

\paragraph{Comparators and additional trajectories.}
The five structured architectures retain the definitions in
\cref{sec:experiment4f}. Two additional clustered LTI baselines fit constant
matrices directly: LTI-all uses all three training amplitudes, reproducing 4D's
estimation approach, while LTI-near uses precisely the same near-linear records
as structured M4. A training-physics M4 transfer controller supplies a reference.
This separates the structured-LTI baseline choice from a claim against the
original directly fitted LTI controllers.

Each frozen controller is evaluated on three 12 s scenarios:
\begin{enumerate}
\item the previous held-out S-curve with $28\rightarrow16\rightarrow24$ m/s;
\item a shifted double lane change with $16\rightarrow27\rightarrow19$ m/s,
using opposite Gaussian yaw pulses of amplitude $8.5^\circ$/s, centered at
3.8 and 7.4 s with widths 0.85 and 0.65 s;
\item an unseen constant speed of 22 m/s and an enveloped sinusoidal yaw
reference, in degrees/s,
\[
 r^\star(t)=7.5\sin^2(\pi t/12)\sin(2\pi\,0.18t+0.35).
\]
\end{enumerate}
Every reference is scaled by 0.5, 1, and 1.5. The labels near-linear, moderate,
and strong denote these inherited amplitude factors. The additional maneuvers
have not been recalibrated to equal tire-force utilization. Speed remains in
the identification support. These tests concern interpolation and new maneuver
combinations, not extrapolation to new physical families or speeds.

\paragraph{Statistical unit and decision rule.}
Moderate-severity tracking is primary. For each train/test seed pair, average
equally over the three scenarios and the balanced test fleet. Paired differences
are baseline RMSE minus M4 RMSE. Confidence intervals resample the ten seed pairs
10,000 times with bootstrap seed 4107, preserving all within-fleet dependencies.
They are pointwise 95\% bootstrap intervals, not simultaneous significance
statements. Individual clients and scenarios are not treated as independent
replicates. Per-scenario contrasts and near-linear/strong outcomes are secondary.

The recovery gate requires at least 5\% reduction in the ratio of mean primary
RMSEs and a positive lower confidence bound against structured M2, structured
M3, LTI-all, and LTI-near. M4 must also be amplitude-feasible in every run and
have spectral radius below one on every held-out frozen audit point. The audit
uses 161 speeds in $[12,28]$ m/s and is not a time-varying LPV certificate.
No failed client or seed is excluded. A vectorized simulator accelerates fleet
evaluation and is tested against the original scalar RK4 simulator and metrics.

\paragraph{Results and interpretation.}
All ten seed pairs completed, yielding 21,600 closed-loop evaluations across
eight methods, three scenarios, three amplitude factors, and 300 held-out
clients. \Cref{tab:experiment4g-means} reports the primary fleet averages and
variation across seed pairs. Structured M4 nearly matches the training-physics
controller, while outperforming both structured alternatives and the original
directly fitted LTI baselines. Its average worst-family error is also smaller
than those of the learned comparators.

\begin{table}[htbp]
\centering
\caption{Experiment 4G: moderate tracking, averaged equally over three scenarios.
Mean, seed standard deviation, and mean worst-family RMSE are in degrees/s.
Worst-family values take the maximum family mean within each seed pair before
averaging across pairs.}
\label{tab:experiment4g-means}
\begin{tabular}{lrrr}
\toprule
Method & Mean $\pm$ seed SD & Worst family & Worst frozen $\rho$\\
\midrule
Structured M1 & $0.03723\pm0.00067$ & 0.04487 & 0.9625\\
Structured M2 & $0.03341\pm0.00066$ & 0.03430 & 0.9653\\
Structured M3 & $0.03222\pm0.00063$ & 0.03766 & 0.9648\\
Structured M4 & $\mathbf{0.02757}\pm0.00041$ & 0.02842 & 0.9668\\
Structured M5 & $0.03342\pm0.00057$ & 0.03455 & 0.9668\\
Direct LTI-all & $0.03528\pm0.00242$ & 0.05430 & 0.9639\\
Direct LTI-near & $0.03314\pm0.00224$ & 0.05044 & 0.9625\\
Physics M4 reference & $0.02752\pm0.00043$ & 0.02830 & 0.9668\\
\bottomrule
\end{tabular}
\end{table}

\Cref{tab:experiment4g-paired} gives the paired evidence. All four primary
contrasts favor M4 in all ten seed pairs, with strictly positive interval lower
bounds and reductions exceeding 5\%. Thus the frozen recovery gate passes.
The direct-LTI comparisons are especially relevant: the improvement is not
confined to the frozen structured LTI baseline introduced in 4F.

\begin{table}[htbp]
\centering
\caption{Experiment 4G paired moderate tracking contrasts. $\Delta$ is baseline
minus M4 RMSE in degrees/s. Intervals bootstrap the ten complete seed pairs.}
\label{tab:experiment4g-paired}
\begin{tabular}{lrrr}
\toprule
Baseline & M4 reduction & Mean $\Delta$ & 95\% interval for $\Delta$\\
\midrule
Structured M2 & 17.5\% & 0.00584 & $[0.00565,0.00602]$\\
Structured M3 & 14.4\% & 0.00465 & $[0.00433,0.00497]$\\
Direct LTI-all & 21.9\% & 0.00771 & $[0.00637,0.00901]$\\
Direct LTI-near & 16.8\% & 0.00558 & $[0.00433,0.00677]$\\
\bottomrule
\end{tabular}
\end{table}

\Cref{fig:experiment4g} shows that the aggregate result is supported by each
scenario. Against structured M2, M4 improves by 9.1\% on the original S-curve,
31.1\% on the shifted lane change, and 9.8\% at constant 22 m/s. Against
direct LTI-near, the reductions are 11.3\%, 28.4\%, and 8.5\%. The last
comparison favors M4 in nine of ten individual pairs, so scenario-wise
dominance is not universal. Constant-speed performance reflects selecting an
appropriate scheduled gain away from the LTI design speed, not time variation
within that particular run. M5 chooses its lower nearest anchor at the 22 m/s
tie, which explains its equality with structured M2 in this scenario.

\begin{figure}[htbp]
\centering
\includegraphics[width=\linewidth]{experiment_4g_validation.pdf}
\caption{Experiment 4G. Left: moderate tracking averaged over scenarios, with
seed standard deviations. Right: relative M4 benefits by scenario. Dashed gray
marks the 5\% aggregate practical threshold, not a separate scenario gate.}
\label{fig:experiment4g}
\end{figure}

All methods remain amplitude-feasible in every evaluated trajectory. Structured
M4 has worst held-out frozen spectral radius 0.9668. Its near-linear and strong
mean RMSEs are $0.01321\pm0.00020$ and $0.04532\pm0.00068$ degrees/s. The
physics reference is approximately 0.17\% better on the aggregate moderate
metric. Consequently, the result supports recovery of useful physics-controller
performance from data, not superiority to known physics. Absolute differences
are small and must be considered alongside their relative percentages.

\paragraph{Scope and next decision.}
The independent-fleet recovery is supported within the original family
distribution, known geometry, noise-free state measurements, two training
profile pairings, and three deterministic test scenarios. It does not establish
robustness to new family centers, sensor noise, limited speed coverage, learned
family assignments, or arbitrary scheduling rates. No formal time-varying
stability certificate is claimed. The structured identification route now has
stronger oracle support than 4F alone. The subsequent complexity and regime
studies can use this validated route, while finite local information and
federated recovery remain distinct questions. No federated learning claim follows
from this centralized oracle experiment.
